\section{CONCLUSIONES}

En las aguas servidas de la PTAR Quitumbe, este estudio obtuvo evidencia molecular y genómica de los tres virus analizados, con un alcance que dependió de la carga viral recuperada en cada caso. La RT-qPCR detectó SARS-CoV-2 en la mayoría de las semanas de 2025, y la secuenciación Oxford Nanopore caracterizó genómicamente seis muestras tempranas, con linajes descendientes de JN.1 en los clados 24A (JN.1, LZ.5, LU.2.1.1) y 24H (LF.7.1.3) \cite{MejiaCalle_2024,Crits_Christoph_2021}. En el caso del enterovirus, el flujo de PCR anidada caracterizó un Coxsackie A, un enterovirus de interés clínico, en dos muestras prospectivas \cite{Bubba_2023,Asghar_2014}. El RSV, en cambio, solo dejó señal molecular en cuatro muestras, con una cobertura insuficiente para asignar el subgrupo \cite{Le_2025,Hughes_2022}. La evidencia genómica fue robusta para SARS-CoV-2 en condiciones de alta carga y se limitó al nivel de detección o de presencia para el enterovirus y el RSV.

\subsection{Limitaciones del estudio}

La principal limitación se debió a la escasa cantidad de material viral recuperado de una matriz ambiental diluida. Cuando los valores de Cq fueron altos, la secuenciación produjo consenso con cobertura insuficiente para asignar linaje, lo que separó la detección molecular de la caracterización genómica en SARS-CoV-2 desde la semana 10 y dejó a RSV sin caracterización \cite{Parkins_2023,Child_2023}. Las muestras retrospectivas conservadas en DNA/RNA Shield 2X y en PBS 1X no rindieron secuencia de poliovirus ni de RSV, lo que indica degradación o carga insuficiente durante el almacenamiento y restringe el uso de estos esquemas para el análisis genómico diferido \cite{MejiaCalle_2024}. El diseño temporal presentó ventanas de análisis con escaso solapamiento entre virus, y SARS-CoV-2 no se analizó entre las semanas 43 y 56, lo que impide evaluar la co-circulación simultánea de los tres patógenos. La disponibilidad limitada de datos genómicos clínicos de SARS-CoV-2 en Ecuador restringió la comparación entre la señal ambiental y la circulación documentada en pacientes \cite{Bruno_2025}. El análisis de enterovirus y RSV se sustentó en 22 muestras de un único sitio centinela, y los sesgos de amplificación propios de los paneles dirigidos pudieron afectar la representación de las variantes de baja abundancia \cite{Child_2023}. Además, en enterovirus los cebadores dirigidos a VP1 amplificaron la región VP4 y no permitieron confirmar el serotipo, lo que evidencia una limitación de especificidad del protocolo en la matriz ambiental.

\subsection{Importancia y recomendaciones}

El estudio aporta la primera línea de base molecular y genómica multipatógeno de aguas servidas en Quito que integra SARS-CoV-2, enterovirus y RSV en una misma matriz ambiental. La detección de los tres virus con un flujo común y la caracterización de un enterovirus de interés clínico confirman la viabilidad operativa de un esquema multipatógeno en un sitio centinela urbano y muestran el valor de la vigilancia ambiental para reconocer circulación que la vigilancia clínica no captura \cite{Bubba_2023,Asghar_2014,Parkins_2023}.

Para estudios futuros se recomienda iniciar con el procesamiento de las muestras de manera temprana frente a la conservación prolongada, dada la pérdida observada en las muestras retrospectivas. La selección de muestras para secuenciación mediante valores de Cq obtenidos por RT-qPCR y la amplificación del genoma en la PCR previa permite concentrar el esfuerzo en las muestras con mayor probabilidad de rendir secuencias útiles \cite{MejiaCalle_2024,Parkins_2023}. La incorporación de plataformas de secuenciación de lectura corta o de estrategias de enriquecimiento del material biológico podría mejorar la recuperación de genomas en muestras de baja carga \cite{Child_2023}. El aumento del número de muestras y de réplicas para enterovirus y RSV, el trabajo en conjunto con datos clínicos y de repositorios genómicos, y el mantenimiento de la vigilancia de enterovirus mediante la región VP1 fortalecerían la capacidad del esquema para sostener una vigilancia genómica ambiental continua \cite{Le_2025,Bruno_2025,Klapsa_2022}.
